\begin{abstract}
    Maritime Search and Rescue operations demand rapid decisions while under significant pressure.
    Due to the severe nature of these tasks, operator workload is more sensitive to increases in task complexity and time pressure, which can impact mission outcomes.
    This paper presents BEACON, a Behavioural Engine for Autonomous Coordination in Ocean Navigation, a digital twin-inspired supervisory framework for coordinating drone swarms in MSAR-like scenarios.
    Swarm Coverage is achieved using a weighted Voronoi partitioning approach, with faster drones assigned proportionally larger areas.
    Within these Voronoi partitions, Boustrophedon path planning is used to generate predictable and explainable autonomous coverage paths.
    These are then displayed as a series of waypoints completed sequentially on the map.
    The interface combines a mission map, a structured 4-tier alert system, and real-time performance evaluation through operational metrics (e.g., detection rate, coverage, alerts/minute).
    These metrics are displayed with an optional NASA-TLX workload assessment to provide a comprehensive view of both system performance and operator workload.

    BEACON seeks to provide a coherent and extensible testbed for studying workload-aware human-swarm supervision in the context of MSAR.
    However, conclusions are limited since validation is currently limited to simulation-based analysis, without human-subject trials or physical drone testing.
\end{abstract}
